% Generated draft: requires mathematical and source review.
\documentclass[11pt]{article}
\usepackage[margin=25mm]{geometry}
\usepackage{fontspec}
\setmainfont{Latin Modern Roman}
\usepackage{amsmath,amssymb,mathtools,unicode-math}
\setmathfont{Latin Modern Math}
\usepackage{enumitem,needspace,xurl,hyperref}
\setlength{\parindent}{0pt}
\setlength{\parskip}{5pt}
\setlength{\emergencystretch}{3em}
\newcommand{\N}{\mathbb{N}}
\newcommand{\Z}{\mathbb{Z}}
\newcommand{\Q}{\mathbb{Q}}
\newcommand{\R}{\mathbb{R}}
\newcommand{\C}{\mathbb{C}}
\newcommand{\F}{\mathbb{F}}
\newcommand{\A}{\mathbb{A}}
\newcommand{\PP}{\mathbb{P}}
\newcommand{\E}{\mathbb{E}}
\newcommand{\eps}{\varepsilon}
\newcommand{\dd}{\,\mathrm{d}}
\begin{document}
\section*{UnsolvedMath ID: 2866}
\textbf{Source problem code:} catalogue entry
\subsection{Definitions and mathematical statement}
The source asks for two constructions: (a) a mathematical definition of the proposed 3-manifold invariants \(\widehat Z\) for all 3-manifolds, and (b) a categorification of those invariants. The background denotes the expected series by \(\widehat Z_a(M;q)\), with integral coefficients and convergence in \(|q|<1\). It expects appropriate limits at certain roots of unity to recover Witten--Reshetikhin--Turaev invariants.

The label set for \(a\), the grading/exponent normalization of the series, and the exact class of manifolds and auxiliary data are not supplied. In particular, ``all 3-manifolds'' does not state closedness, compactness, orientation, boundary conditions or the treatment of complements. No unmentioned spin, framing or boundary decoration is invented to turn that phrase into a narrower theorem.

A precise all-manifold definition would need to specify the admissible labelled objects, the series associated to each, and its invariance under the intended equivalence of labelled manifolds. For a specified root \(\zeta\), a radial limit means approach through \(q=\rho\zeta\), \(0<\rho<1\), as \(\rho\uparrow1\). The source does not state the exact roots, required prefactors or combinations of labelled series that recover a particular WRT normalization. Consequently it does not justify an invented termwise equality of every individual raw radial limit with a WRT invariant.

Categorification must likewise specify the category, gradings and decategorification map. One possible homological interpretation would attach graded groups \(H^{i,j}(M,a)\), with finiteness sufficient to define
\[
\sum_j\left(\sum_i(-1)^i\operatorname{rank}H^{i,j}(M,a)\right)q^j,
\]
and identify this with the normalized \(\widehat Z_a\). This is an illustrative meaning of categorification, not the only categorical model or an additional theorem imposed by the source. The background's integrality is motivation for such a construction, not itself an invariant homology theory, nor a request for arbitrary coefficient-matching vector spaces without topological meaning.

The supplied triage dated 2026-08-17 reports rigorous definitions in negative-definite plumbing settings, Gukov--Manolescu's two-variable series for knot complements, selected surgeries and mapping-torus contexts. It separately reports Sugimoto's 2025 Seifert and later negative-definite-plumbing categorical blueprints as proposals/preprints. These restricted constructions are not promoted to definitions and categorifications for every 3-manifold. The background's large-\(N\) connection to enumerative geometry is a separate suggested direction without a parameter definition in the question. The registry status is open; the dated triage labels the row partially solved and leaves both all-manifold tasks open, including full invariance and WRT radial-limit proofs. No categorical proposal or current status is independently certified here.
\subsection{Short English statement}
Define the conjectural integral \(q\)-series \(\widehat Z\) for all intended 3-manifolds, prove its invariance and specified WRT-limit relation, and construct a topological categorification. The source reports rigorous restricted-family formulas and 2025 categorical proposals, while leaving the manifold class, labels, normalization and full general construction unresolved.
\subsection{Sources}
\begin{itemize}
\item Local UnsolvedMath registry, permanent ID 2866, source code catalogue entry, ``Kirby Problem 3.68,'' including supplied literature triage dated 2026-08-17.
\item Baykur, Kirby, and Ruberman, \emph{K3: A New Problem List in Low-Dimensional Topology}, Problem 3.68 (2026): \url{https://bpb-us-e2.wpmucdn.com/websites.umass.edu/dist/b/22144/files/2026/04/K3-problem-list-watermarked.pdf}.
\item Sergei Gukov, Pavel Putrov, and Cumrun Vafa, ``Fivebranes and 3-manifold homology,'' \emph{Journal of High Energy Physics} 2017(7), 071, front matter plus 80 pp. (2017), DOI 10.1007/JHEP07(2017)071, source label [GPV17].
\item Sergei Gukov, Du Pei, Pavel Putrov, and Cumrun Vafa, ``BPS spectra and 3-manifold invariants,'' \emph{Journal of Knot Theory and Its Ramifications} 29(2) (2020), 2040003, 85 pp., DOI 10.1142/S0218216520400039, source label [GPPV20].
\item Sergei Gukov and Ciprian Manolescu, ``A two-variable series for knot complements,'' \emph{Quantum Topology} 12(1) (2021), 1--109, DOI 10.4171/qt/145, source label [GM21]: \url{https://doi.org/10.4171/QT/145}.
\item Sungbong Chun, Sergei Gukov, Sunghyuk Park, and Nikita Sopenko, ``3d-3d correspondence for mapping tori,'' \emph{Journal of High Energy Physics} 2020(9), 152, 59 pp. (2020), DOI 10.1007/jhep09(2020)152, source label [CGPS20].
\item Tobias Ekholm, Angus Gruen, Sergei Gukov, Piotr Kucharski, Sunghyuk Park, Marko Sto\v{s}i\'c, and Piotr Sul kowski (as listed), ``Branches, quivers, and ideals for knot complements,'' \emph{Journal of Geometry and Physics} 177 (2022), Paper 104520, 75 pp., DOI 10.1016/j.geomphys.2022.104520, source label [EGG+22].
\item Shoma Sugimoto, ``Hypercubic structures behind hat-Z-invariants,'' arXiv:2501.12985 (2025): \url{https://arxiv.org/abs/2501.12985}.
\item Shoma Sugimoto, ``Nesting behind hat-Z-invariants,'' arXiv:2507.13996 (2025): \url{https://arxiv.org/abs/2507.13996}.
\end{itemize}
\end{document}
